\section{Low-dimensional proofs}
\label{low:appendix}

This appendix supplies the rigidity and curved-path proofs for Sec.~\ref{low:section}, followed by the invariant that distinguishes the specified four-qubit reference measurement from the known larger-radius segment.

\subsection{An exact rigidity certificate at \texorpdfstring{$n=3$}{n=3}}
\label{low:app-rigidity}

Write $b_z=i^{f_0(z)}$ and let $\theta_0=0$.  For a signed physical
Pauli operator $\widetilde P$, define the fixed matrix coefficients
\begin{equation}
 K_P(z,w)=\overline{b_w}b_z
 \langle w|C_n^\dagger\widetilde P C_n|z\rangle.
 \label{low:finite-coeff}
\end{equation}
The seed component and all its derivatives follow from the finite identity
\begin{equation}
 2^n m_P(\theta)=\sum_{z,w}K_P(z,w)e^{i(\theta_z-\theta_w)}.
 \label{low:finite-sum}
\end{equation}
In particular,
\begin{equation}
 2^n\partial_lm_P(0)
 =i\sum_{z,w}K_P(z,w)
 (\mathbf1_{z=l}-\mathbf1_{w=l}).
 \label{low:jacobian-formula}
\end{equation}
These formulas use the original circuit and therefore keep its signs and
normalization explicit.

For $n=3$, let $\mathcal F$ contain the eight differences $m_P-m_{X_1}$,
in the order
\begin{equation*}
 P=Y_1,Z_1,X_2,Y_2,Z_2,X_3,-Y_3,Z_3.
\end{equation*}
The columns correspond to $\theta_1,\ldots,\theta_7$.
Equation~\eqref{low:jacobian-formula} gives
\begin{equation}
 4D\mathcal F(0)=
 \begin{pmatrix}
 -2&-2&0&1&1&0&2\\
 -1&-2&0&-1&2&2&0\\
 -1&-1&-1&0&1&0&1\\
 -1&-1&-1&0&1&1&0\\
 -1&-2&0&1&0&0&2\\
 -2&-1&0&1&0&0&2\\
 0&-2&-1&1&2&0&0\\
 -1&0&-2&1&0&2&0
 \end{pmatrix}.
 \label{low:jacobian-three}
\end{equation}
The determinant of the first seven rows of the displayed integer matrix is
$4$.  Hence the corresponding derivative determinant is
$4/4^7=1/4096$.  The inverse function theorem applies to these seven
real analytic component differences and implies that their common zero is
locally unique.  Every regular-tetrahedral measurement sufficiently near
the specified point satisfies these equations: all signed components are
positive there, so equality of their squares implies equality of the
components themselves.  This proves Proposition~\ref{low:rigidity} in
the full seven-dimensional phase space.

\subsection{All real phase lines at \texorpdfstring{$n=4$}{n=4}}
\label{low:app-lines}

Let $\mathcal F$ now contain the eleven differences from $X_1$, with the
remaining signed operators ordered as
\begin{equation*}
 Y_1,Z_1,X_2,Y_2,Z_2,X_3,Y_3,Z_3,X_4,-Y_4,Z_4.
\end{equation*}
For a line $\theta=td$, Eq.~\eqref{low:finite-sum} gives the exact
derivative rule
\begin{equation}
 16\frac{d^k}{dt^k}m_P(td)\bigg|_{t=0}
 =\sum_{z,w}K_P(z,w)\bigl[i(d_z-d_w)\bigr]^k.
 \label{low:line-derivative}
\end{equation}
The first-order equations have precisely the following four-parameter
solution, where $a,b_*,c,\kappa$ are arbitrary real numbers:
\begin{equation}
 \begin{aligned}
 d_0&=0,\qquad d_1=d_2=d_9=a,\\
 d_3=d_8&=a-c,\qquad d_{10}=b_*-c,\\
 d_4=d_5&=-a+b_*-c/2+\kappa,\\
 d_6=d_7=d_{12}=d_{13}&=c,\\
 d_{11}&=b_*,\qquad d_{14}=d_{15}=\kappa.
 \end{aligned}
 \label{low:kernel-four}
\end{equation}
For completeness, this is an exact linear-algebra reduction of
Eq.~\eqref{low:jacobian-formula}: all four displayed directions annihilate
$D\mathcal F(0)$, and the minor of $16D\mathcal F(0)$ with columns
$(1,2,3,4,5,6,7,8,10,12,14)$ has determinant $65536$.
Thus its rank is $11$ and the displayed kernel is exhaustive.

Put $G_P(t)=16[m_P(td)-m_{X_1}(td)]$.  The second-order conditions from
Eq.~\eqref{low:line-derivative} include
\begin{align}
 G''_{Y_1}(0)&=-2c(2a-c),\nonumber\\
 G''_{Y_2}(0)&=4c(a-2b_*+2c).
 \label{low:second-order}
\end{align}
If $c\ne0$, their vanishing requires $a=c/2$ and $b_*=5c/4$.
The next condition becomes
\begin{equation}
 G''_{X_2}(0)=-\frac{c(c+4\kappa)}2=0,
\end{equation}
so $\kappa=-c/4$.  But the same derivative rule then gives
\begin{equation}
 G'''_{Y_1}(0)=-\frac{39c^3}{8}\ne0.
 \label{low:third-order}
\end{equation}
Consequently $c=0$.  The remaining necessary second-order condition is
\begin{equation}
 G''_{X_2}(0)=-8(a-b_*)(a-b_*-\kappa)=0.
 \label{low:two-planes}
\end{equation}
Define the three subsets of the integer input labels
\begin{align}
 \mathcal A&=\{1,2,3,8,9,10,11\},\nonumber\\
 \mathcal B&=\{1,2,3,8,9,14,15\},\nonumber\\
 \mathcal C&=\{4,5,14,15\}.
 \label{low:line-supports}
\end{align}
The branch $b_*=a$ gives $d=a\mathbf1_{\mathcal A}
+\kappa\mathbf1_{\mathcal C}$, whereas the branch
$\kappa=a-b_*$ gives $d=b_*\mathbf1_{\mathcal A}
+(a-b_*)\mathbf1_{\mathcal B}$.  Substitution into the finite sum
\eqref{low:finite-sum} gives, for every signed local Pauli operator and
either branch,
\begin{equation}
 16m_P(td)=e^{iat}+e^{-iat}.
 \label{low:planes-sufficient}
\end{equation}
This checks sufficiency for all real parameters, rather than merely to
finite order.  Hence all allowable directions are exactly
\begin{equation}
 \operatorname{span}_{\mathbb R}
 \{\mathbf1_{\mathcal A},\mathbf1_{\mathcal C}\}
 \ \cup\ 
 \operatorname{span}_{\mathbb R}
 \{\mathbf1_{\mathcal A},\mathbf1_{\mathcal B}\}.
 \label{low:line-classification}
\end{equation}
Local regularity implies the signed equations used above by positivity at
the base point.  Their differences along any real direction are finite
exponential sums, so equality on an open interval extends to the whole
line.  Equations~\eqref{low:planes-sufficient} and
\eqref{low:line-classification} establish the first part of
Theorem~\ref{low:theorem} without a restriction on the number of values
taken by $d$.

\subsection{The curved branch and its full monotonic interval}
\label{low:app-curve}

Equation~\eqref{low:masks} specifies all phases in Eq.~\eqref{low:curve}.
Its entries are indexed by $z=0,\ldots,15$, with the most significant
bit first.  Inserting $\theta_z=bu_z+\delta v_z$ into
Eq.~\eqref{low:finite-sum} gives precisely the two expressions
in Eq.~\eqref{low:two-components}.  In particular, these identities
hold for independent real $b,\delta$, before any branch is selected.

\begin{equation}
 \begin{array}{c|rrrrrrrr}
 z&0&1&2&3&4&5&6&7\\
 u_z&0&1&1&1&2&2&0&0\\
 v_z&0&1&1&0&1&1&1&1\\[2pt]
 z&8&9&10&11&12&13&14&15\\
 u_z&1&1&0&0&0&0&1&1\\
 v_z&0&1&0&1&1&1&0&0
 \end{array}
 \label{low:masks}
\end{equation}

Set $q=-\delta$ and $z=\tan(q/2)$; the letter $z$ denotes a real
variable only in the following branch calculation.  The half-angle
identities give
\begin{equation}
 H-T=\frac{2(Az^2-Bz+C)}{1+z^2}.
 \label{low:quadratic-identity}
\end{equation}
We now verify the root in Eq.~\eqref{low:delta}, including its entire
domain.  Put $x=\tan(b/2)$ and $k=\sqrt2-1$.  Then
\begin{equation}
 Az^2-Bz+C=\frac{x+1}{(1+x^2)^2}F(x,z),
 \label{low:halfangle-polynomial}
\end{equation}
where
\begin{align}
 F(x,z)={}&-x^3z^2+x^3z-3x^2z^2-3x^2z-4x^2\nonumber\\
 &+3xz^2+xz+4x+z^2-3z.
 \label{low:polynomial}
\end{align}
On the square $0\le x,z\le k$, one has
\begin{equation}
 F_{xx}=6z[x(1-z)-(1+z)]-8<0.
\end{equation}
Using $k^2+2k-1=0$ and $4-12k<0$, it follows that
\begin{align}
 F_x(x,z)&\ge F_x(k,z)
 =4-8k+(4-12k)z\nonumber\\
 &\ge20k-8=20\sqrt2-28>0.
 \label{low:fx-positive}
\end{align}
The other derivative satisfies
\begin{align}
 F_z&=2z(1+3x-3x^2-x^3)+x^3-3x^2+x-3\nonumber\\
 &\le2k(1+3k)+k^3+k-3
 =5-4\sqrt2<0.
 \label{low:fz-negative}
\end{align}
The boundary values obey
\begin{align}
 F(x,x)&=-x(1+x^2)(x^2+2x-1)\ge0,\nonumber\\
 F(x,k)&\le F(k,k)=0.
 \label{low:boundary}
\end{align}
Thus each $x\in[0,k]$ has a unique root $z(x)\in[x,k]$;
$z(0)=0$, $z(k)=k$, and $x<z(x)<k$ in the interior.
At the root, $F_z<0$.  Since $A>0$, the corresponding quadratic has
positive discriminant and this is its smaller root.  Rationalizing that
root gives
\begin{equation}
 z(x)=\frac{2C}{B+\sqrt{B^2-4AC}},
 \label{low:root-identified}
\end{equation}
which is well defined also at $x=0$.  This proves the stated closed
formula rather than defining the path by an implicit existence argument.

Differentiating the now-identified simple root gives
$z'(x)=-F_x/F_z>0$.  Consequently $q'(b)>0$ and
$b\le q(b)\le\pi/4$.  With $T=\cos(b-q)+\sin q$,
\begin{equation}
 T'(b)=\sin(q-b)+q'(b)\,[\cos q-\sin(q-b)]>0.
 \label{low:monotone}
\end{equation}
The bracket is strictly positive: equality would require both
$q=\pi/4$ and $b=0$, inconsistent with the determined endpoint values.
The first term is nonnegative.  Finally,
$T(0)=1$ and $T(\pi/4)=1+1/\sqrt2$ prove the exact image
\eqref{low:curve-range}.  At the origin,
$z'(0)=4/3$, $q'(0)=4/3$, and $r'_4(0)=\sqrt3/6$.

Setting $n=4$ and $m=T/8$ in Eq.~\eqref{core:bloch-orbit} gives all
reduced states, with each fixed tetrahedral direction occurring four times.
The unitary factorization in Eq.~\eqref{low:curve} gives a complete
orthonormal measurement, and $b=0$ recovers the specified $D_{f_0}$.

\subsection{Non-equivalence of the reference measurement to the known \texorpdfstring{$n=4$}{n=4} segment}
\label{low:app-invariant}

We show that the specified reference measurement is not equivalent, under local unitaries, qubit permutations, or outcome relabeling, to any measurement on the $n=4$ segment in Appendix C of Ref.~\cite{PauwelsGisin2026}.
With $a\in[0,\pi/3]$, let $x=e^{ia}$ and write its input phases as
$i^{f_z}x^{k_z}$.  In order $z=0,\ldots,15$, the two fixed vectors are
\begin{align}
 f={}&(0,0,0,3,0,0,1,0,\nonumber\\
 &\hspace{12mm}0,0,0,2,0,1,3,3),\nonumber\\
 k={}&(0,1,-1,-1,0,1,0,0,\nonumber\\
 &\hspace{12mm}-1,2,0,1,3,3,1,0).
 \label{low:f159-data}
\end{align}
Here $a=\pi(1-\tau)/3$ in the parameterization of that reference.
The seed $C_4(i^{f_z}x^{k_z})_z/4$ has one-qubit eigenvalues
$(1\pm\sqrt3\,c(x))/2$, where
\begin{equation}
 c(x)=\frac{x^3+x^{-3}+2x+2x^{-1}}{16}
 =\frac{\cos3a+2\cos a}{8}.
 \label{low:f159-component}
\end{equation}
This is nonnegative on the stated segment, decreases from $3/8$ to
zero, and gives the already known endpoint $3r_E(4)$.

For a pure state at $n=4$ define
\begin{equation}
 S(\psi)=\operatorname{Tr}(\rho_{12}^{2})
 +\operatorname{Tr}(\rho_{13}^{2})
 +\operatorname{Tr}(\rho_{14}^{2}).
 \label{low:purity-invariant}
\end{equation}
The sum is invariant under local unitaries and qubit permutations:
complementary subsystems of a pure state have identical nonzero spectra,
so a permutation only permutes the three unordered $2|2$ terms.
Each column of a Pauli-orbit measurement shares these invariants with
its seed.  Outcome relabeling therefore cannot change the comparison.
At the specified $f_0$ point the three purities are
$(63,39,39)/128$, and hence $S_0=141/128$.

An equivalent measurement must have both the same one-qubit spectrum
and the same $S$.  On this segment these conditions require $c=1/8$
and $S=141/128$, respectively. Evaluating the reduced-state purities
from the amplitudes in Eq.~\eqref{low:f159-data} and clearing
denominators, with $x\ne0$, gives the necessary equations
$P(x)=Q(x)=0$, with
\begin{align}
 P(x)={}&x^6+2x^4-2x^3+2x^2+1,\nonumber\\
 Q(x)={}&6x^{16}-3x^{14}+74x^{12}+48x^{10}\nonumber\\
 &+134x^8+48x^6+74x^4-3x^2+6.
 \label{low:invariant-polynomials}
\end{align}
The Euclidean algorithm excludes a common root. Its first remainder is
\begin{equation}
 \begin{aligned}
 Q\bmod P={}&816x^5-2340x^4+2168x^3\\
 &-1240x^2+616x-596,
 \end{aligned}
\end{equation}
whose gcd with $P$ is $1$.  Thus the original parameter segment
contains no measurement equivalent to the specified $f_0$ measurement
under the allowed operations.





